% Compile with pdfLaTeX on Overleaf.
% Report for: Phan tich cam xuc theo khia canh va sinh bao cao
% tong hop tu phan hoi khach hang bang Generative AI (20252B - Gen AI)
%
% TODO markers throughout mark content owned by the report-generation /
% factual-checker team (Hoang, Vinh, Hung) or not-yet-done analysis --
% fill in once notebooks-output/report_generation_semeval_laptop_output.ipynb
% and the factual-checker branches are merged, and after error analysis.
%
% All numeric results/figures below are taken directly from
% plans/project-plan.md and the executed notebooks in notebooks-output/ --
% no fabricated numbers. Charts are drawn with pgfplots directly from the
% same logged per-epoch metrics, no external image files needed.

\documentclass[11pt,a4paper]{article}

\usepackage[utf8]{inputenc}
\usepackage[T5]{fontenc}
\usepackage[vietnamese]{babel}
\usepackage[margin=2.5cm]{geometry}
\usepackage{graphicx}
\usepackage{booktabs}
\usepackage{amsmath}
\usepackage{hyperref}
\usepackage{xcolor}
\usepackage{listings}
\usepackage{caption}
\usepackage{tikz}
\usepackage{pgfplots}
\pgfplotsset{compat=1.17}
\usetikzlibrary{arrows.meta}

\newcommand{\todo}[1]{\textcolor{red}{\textbf{[TODO: #1]}}}

\title{Phân tích cảm xúc theo khía cạnh và sinh báo cáo tổng hợp \\
từ phản hồi khách hàng bằng Generative AI}
\author{Nhóm 20252B \\ \small Đạt, Sơn, Hoàng, Vinh, Hưng}
\date{\today}

\begin{document}

\maketitle

\begin{abstract}
Thay vì chỉ phân loại một review là tích cực/tiêu cực, hệ thống trong báo cáo này
xác định cụ thể khách hàng đang khen/chê điểm nào (khía cạnh -- \emph{aspect}),
gán sentiment cho từng khía cạnh, rồi tổng hợp thống kê trên nhiều review và dùng
mô hình sinh văn bản để tạo báo cáo tóm tắt ngắn, kèm một bộ kiểm tra tính xác
thực (factual checker) để hạn chế mô hình bịa số liệu. Nhóm xây dựng và so sánh
hai hướng tiếp cận trên hai domain: (1) phân loại sentiment theo khía cạnh đã
biết trước (TF-IDF+LogReg $\to$ DistilBERT/BERT) trên domain laptop
(SemEval-2014 Task~4), và (2) trích xuất trực tiếp bộ ba (khía cạnh, ý kiến,
sentiment) từ câu thô không cần biết trước khía cạnh (họ mô hình T5) trên domain
nhà hàng (SemEval Restaurant Triplet data). Hướng thứ hai giải quyết luôn bài
toán trích xuất khía cạnh và cho thêm ``lý do'' (opinion phrase) phục vụ báo
cáo tốt hơn, đồng thời cho độ nhất quán tổng hợp cao hơn (97.25\% so với
93.28\%). Nhóm cũng ghi lại một số phát hiện đáng chú ý trong quá trình thực
nghiệm, trong đó có một kết quả phủ định giả thuyết ban đầu về nguyên nhân
FLAN-T5-base hoạt động kém trên tác vụ trích xuất.
\end{abstract}

\tableofcontents

% =====================================================================
\section{Giới thiệu}
% =====================================================================

\subsection{Chủ đề và động lực}

Các hệ thống phân tích cảm xúc truyền thống thường chỉ đưa ra một nhãn
tích cực/tiêu cực duy nhất cho toàn bộ review, trong khi một review thực tế
thường đề cập nhiều khía cạnh khác nhau với sentiment khác nhau. Với người bán
hàng hay quản lý nhà hàng, một con số tổng quát (``70\% review tích cực'')
không đủ để biết \emph{nên cải thiện điểm gì} -- họ cần biết cụ thể khía cạnh
nào đang được khen, khía cạnh nào đang bị chê, và vì sao. Đây chính là động lực
của bài toán \emph{Aspect-Based Sentiment Analysis} (ABSA), minh hoạ qua ví dụ:

\begin{quote}
\emph{"Máy có màn hình đẹp, giao hàng nhanh nhưng pin yếu."}
\end{quote}

\begin{table}[h]
\centering
\begin{tabular}{ll}
\toprule
Khía cạnh & Sentiment \\
\midrule
Màn hình & Tích cực \\
Giao hàng & Tích cực \\
Pin & Tiêu cực \\
\bottomrule
\end{tabular}
\caption{Ví dụ phân tích cảm xúc theo khía cạnh (Aspect-Based Sentiment Analysis).}
\end{table}

Khi có hàng trăm, hàng nghìn review như vậy, việc đọc thủ công từng câu là
không khả thi. Nhóm hướng tới một hệ thống tự động hoá toàn bộ quy trình: từ
review thô, tự tổng hợp và sinh ra một đoạn báo cáo dạng:

\begin{quote}
\emph{"Khách hàng đánh giá tích cực về màn hình và tốc độ giao hàng. Tuy
nhiên, pin là khía cạnh có nhiều phản hồi tiêu cực nhất và cần được ưu
tiên cải thiện."}
\end{quote}

Một rủi ro lớn khi dùng mô hình sinh văn bản (generative model) cho bước cuối
là hiện tượng \emph{hallucination} -- mô hình có thể bịa ra số liệu hoặc nhận
định không khớp với dữ liệu gốc. Vì vậy pipeline của nhóm có thêm một bước
kiểm tra tính xác thực (factual checker) trước khi xuất báo cáo cuối cùng.

\subsection{Pipeline tổng quan}

Hình~\ref{fig:pipeline} minh hoạ toàn bộ pipeline. Từ review thô, nhóm đi theo
hai nhánh xử lý song song (Track~1 và Track~2, xem Mục~\ref{sec:problem}),
kết quả của cả hai nhánh đều đổ về cùng một bước tổng hợp thống kê, sau đó
sinh báo cáo và kiểm tra tính xác thực.

\begin{figure}[h]
\centering
\begin{tikzpicture}[
  box/.style={draw, rounded corners, align=center, minimum width=4.6cm,
    minimum height=1.1cm, font=\small, fill=blue!5},
  chkbox/.style={draw, rounded corners, align=center, minimum width=4.6cm,
    minimum height=1.1cm, font=\small, fill=orange!10, dashed},
  arr/.style={-{Stealth[length=2.2mm]}, thick}
]
\node[box] (input) at (2.3,7.6) {Review khách hàng (văn bản thô)};
\node[box] (t1) at (-2.6,5.6) {\textbf{Track 1}: TF-IDF/BERT\\(aspect biết trước) $\to$ sentiment};
\node[box] (t2) at (7.2,5.6) {\textbf{Track 2}: T5-ASTE\\(câu thô) $\to$ (aspect, opinion, sentiment)};
\node[box] (agg) at (2.3,3.6) {Tổng hợp thống kê theo aspect\\+ top-10 lý do phổ biến nhất};
\node[box] (rep) at (2.3,1.6) {FLAN-T5: sinh báo cáo văn bản};
\node[chkbox] (chk) at (2.3,-0.4) {Factual checker\\(đối chiếu số liệu với thống kê gốc)};
\draw[arr] (input) -- (t1);
\draw[arr] (input) -- (t2);
\draw[arr] (t1) -- (agg);
\draw[arr] (t2) -- (agg);
\draw[arr] (agg) -- (rep);
\draw[arr] (rep) -- (chk);
\end{tikzpicture}
\caption{Pipeline tổng quan. Track~1 (domain laptop) cần biết trước aspect;
Track~2 (domain nhà hàng, hướng bổ sung) tự trích xuất luôn aspect và cho
thêm trường \emph{opinion} làm ``lý do''. Khối nét đứt (factual checker) là
phần do nhóm Hoàng/Vinh/Hưng phụ trách.}
\label{fig:pipeline}
\end{figure}

\subsection{Tổng quan các phương pháp}

Nhóm xây dựng và so sánh nhiều phương pháp theo từng bước của pipeline,
được trình bày chi tiết ở Mục~\ref{sec:methods}:

\begin{itemize}
    \item \textbf{Baseline TF-IDF + Logistic Regression} -- mốc so sánh
        đơn giản, không cần GPU (domain laptop).
    \item \textbf{DistilBERT / BERT fine-tuned} -- phân loại sentiment
        theo khía cạnh đã biết trước, hiểu ngữ cảnh tốt hơn TF-IDF
        (domain laptop).
    \item \textbf{Baseline tra cứu tần suất (AsteLookupBaseline)} -- mốc
        so sánh phi-neural cho bài toán trích xuất bộ ba (domain nhà hàng).
    \item \textbf{T5-small / T5-base / FLAN-T5-base fine-tuned cho ASTE}
        (Aspect Sentiment Triplet Extraction) -- sinh trực tiếp bộ ba
        (khía cạnh, ý kiến, sentiment) từ câu thô, không cần biết trước
        khía cạnh (domain nhà hàng).
    \item \textbf{FLAN-T5 sinh báo cáo} từ bảng thống kê khía cạnh.
        \todo{Hoàng/Vinh/Hưng -- bổ sung chi tiết method sau khi merge
        branch report-generation.}
    \item \textbf{Factual checker} đối chiếu số liệu trong báo cáo sinh
        ra với bảng thống kê gốc.
        \todo{Hoàng/Vinh/Hưng -- bổ sung chi tiết method + kết quả đánh
        giá sau khi merge.}
\end{itemize}


% =====================================================================
\section{Công trình liên quan}
% =====================================================================

\subsection{Aspect-Based Sentiment Analysis}

ABSA thường được chia thành các bài toán con: \emph{Aspect Term Extraction}
(tìm cụm từ chỉ khía cạnh trong câu), \emph{Aspect Term Polarity} (phân loại
sentiment cho từng aspect term), \emph{Aspect Category Detection} và
\emph{Aspect Category Polarity} (với các domain có category cố định, ví dụ
FOOD/SERVICE/PRICE cho nhà hàng)~\cite{semeval2014}. Domain Laptop trong
SemEval-2014 chỉ gán nhãn cho hai bài toán con đầu (aspect term + polarity,
không có category cố định) -- phù hợp với hướng ``trích xuất khía cạnh tự
do'' mà nhóm theo đuổi thay vì phân loại vào một tập category định sẵn.

\subsection{Repo nền: BERT-E2E-ABSA}

Nhóm chọn \textbf{lixin4ever/BERT-E2E-ABSA}~\cite{berte2eabsa} làm repo nền
để kế thừa ý tưởng. Repo này thực hiện ABSA \emph{end-to-end}: một model
BERT duy nhất vừa trích xuất aspect term vừa phân loại sentiment, dùng
unified tagging scheme (BIEOS/BIO/OT) thay vì tách thành hai model rời rạc,
và có sẵn dataset \texttt{laptop14} (chính là SemEval-2014 Laptop).

Hai lựa chọn khác từng được khảo sát nhưng không chọn:
\begin{itemize}
    \item \texttt{songyouwei/ABSA-PyTorch} -- chỉ phân loại sentiment khi
        \emph{đã biết trước} aspect term, không tự trích xuất được, phải
        ghép thêm một model riêng cho bước trích xuất.
    \item \texttt{ScalaConsultants/Aspect-Based-Sentiment-Analysis} -- thư
        viện đóng gói sẵn, khó chỉnh sửa sâu để thể hiện phần cải tiến của
        nhóm.
\end{itemize}

\textbf{Điểm cải tiến của nhóm so với repo nền}: (1) thêm bước tổng hợp
thống kê aspect-sentiment trên \emph{nhiều} review thay vì chỉ dừng ở một
câu; (2) thêm module sinh report bằng FLAN-T5 từ bảng thống kê; (3) thêm một
factual checker để kiểm tra report có bịa số liệu hay không -- cả ba bước
này đều không có trong repo nền.

\subsection{Aspect Sentiment Triplet Extraction (ASTE)}

Ngoài hướng phân loại sentiment khi biết trước aspect, nhóm còn tìm hiểu
hướng \emph{Aspect Sentiment Triplet Extraction} -- sinh đồng thời bộ ba
(aspect, opinion, sentiment) từ câu thô, không cần biết trước aspect. Nhóm
dùng bộ dữ liệu SemEval Restaurant Triplet
data~\cite{peng2020aste,xu2020aste}, cụ thể phiên bản
\texttt{ASTE-Data-V1-AAAI2020}, với cách gán nhãn theo nhóm token (mỗi
aspect gán một mã nhóm, opinion tương ứng gán mã nhóm cùng độ dài dạng
\texttt{S...S}) để một câu có thể chứa nhiều bộ ba độc lập. Cách tiếp cận
sinh văn bản trực tiếp (thay vì tagging từng token) cho phép tái sử dụng
kiến trúc encoder-decoder có sẵn (T5) mà không cần thiết kế lại tầng CRF/
tagging scheme như repo nền.

\subsection{Mô hình sinh văn bản: T5 và FLAN-T5}

\textbf{T5}~\cite{t5} (Text-to-Text Transfer Transformer) đóng khung mọi
tác vụ NLP dưới dạng text-to-text, phù hợp tự nhiên với bài toán sinh bộ ba
có cấu trúc từ câu thô (Mục~\ref{sec:methods}). \textbf{FLAN-T5}~\cite{flant5}
là phiên bản T5 đã được instruction-tuned trên nhiều tác vụ đa dạng, đã cho
kết quả tốt trên các benchmark table-to-text/data-to-text như
ToTTo~\cite{t5totto} mà không cần fine-tune sâu -- đây là lý do nhóm chọn
FLAN-T5 (thay vì T5 gốc) cho bước sinh báo cáo cuối, dù kết quả thực nghiệm
lại cho thấy FLAN-T5-base \emph{không} phải lựa chọn tốt hơn T5-base cho
bước trích xuất bộ ba (Mục~\ref{sec:findings}) -- một phát hiện thú vị được
thảo luận riêng.

\subsection{Kiểm tra tính xác thực (factual checking)}

Trong literature có nhiều hướng phức tạp để kiểm tra tính nhất quán của văn
bản sinh ra so với nguồn dữ liệu gốc, ví dụ FactCC~\cite{factcc} hay các
phương pháp dùng model riêng để đối chiếu (ConFactCheck và các hướng liên
quan~\cite{consistency2024}). Các hướng này vượt quá phạm vi môn học, nên
nhóm chọn cách tiếp cận rule-based đơn giản hơn: trích số liệu (\%, số
lượng) và tên aspect xuất hiện trong report bằng regex, đối chiếu trực tiếp
với bảng thống kê gốc.
\todo{Hoàng/Vinh/Hưng -- bổ sung mô tả cụ thể cách làm thực tế sau khi
merge, và so sánh ngắn với các hướng phức tạp hơn nêu trên.}


% =====================================================================
\section{Bài toán}
\label{sec:problem}
% =====================================================================

Nhóm giải quyết bài toán trên hai domain song song, với hai định dạng
input/output khác nhau tương ứng hai hướng tiếp cận.

\subsection{Track 1 -- Phân loại sentiment theo khía cạnh (domain laptop)}

\begin{itemize}
    \item \textbf{Input}: một câu review (tiếng Anh) \emph{và} một khía
        cạnh (aspect term) đã biết trước, nằm trong câu đó.
    \item \textbf{Output}: nhãn sentiment cho khía cạnh đó -- một trong
        \{\texttt{positive}, \texttt{negative}, \texttt{neutral}\}.
    \item \textbf{Dataset}: SemEval-2014 Task~4 (Laptop reviews), có sẵn
        nhãn aspect term + polarity gốc trong dữ liệu.
\end{itemize}

\subsection{Track 2 -- Trích xuất bộ ba khía cạnh--ý kiến--sentiment (domain nhà hàng)}

\begin{itemize}
    \item \textbf{Input}: một câu review thô, \emph{không} có thông tin
        khía cạnh đi kèm.
    \item \textbf{Output}: danh sách bộ ba
        (\emph{aspect}, \emph{opinion}, \emph{sentiment}), ví dụ:
        \begin{quote}
        \texttt{aspect: price | opinion: reasonable | sentiment: positive}
        \end{quote}
        Trường \emph{opinion} chính là cụm từ diễn đạt lý do (reason)
        cho sentiment đó, phục vụ trực tiếp cho bước sinh báo cáo.
    \item \textbf{Dataset}: SemEval Restaurant Triplet data (14res +
        15res + 16res).
\end{itemize}

\subsection{Bước tổng hợp và sinh báo cáo (chung cho cả 2 track)}

\begin{itemize}
    \item \textbf{Input}: danh sách (aspect, sentiment[, opinion]) trên
        toàn bộ tập review.
    \item \textbf{Output}: (1) bảng thống kê theo từng aspect (số lượt
        positive/negative/neutral, top-10 lý do phổ biến nhất mỗi
        sentiment), (2) một đoạn báo cáo văn bản ngắn tóm tắt xu hướng,
        (3) kết quả kiểm tra tính xác thực của báo cáo đó.
\end{itemize}


% =====================================================================
\section{Phương pháp}
\label{sec:methods}
% =====================================================================

\subsection{Baseline -- TF-IDF + Logistic Regression (domain laptop)}

Baseline đơn giản nhất được dùng làm mốc so sánh, không cần GPU, giúp định
lượng mức độ cải thiện khi chuyển sang model hiểu ngữ cảnh (BERT). Câu
review được đánh dấu vị trí aspect term bằng token \texttt{\$T\$}, sau đó
vector hoá bằng TF-IDF (n-gram 1--2) trên cả câu lẫn riêng aspect term rồi
ghép hai vector lại làm đặc trưng đầu vào. Bộ phân loại là Logistic
Regression đa lớp của scikit-learn, dùng
\texttt{class\_weight="balanced"} để bù lệch lớp (positive/negative chiếm
đa số so với neutral). Tại thời điểm dự đoán, một vector TF-IDF cho ra xác
suất từng lớp sentiment.

Hạn chế lớn nhất của cách tiếp cận này là không tự trích xuất được aspect
term (bắt buộc phải biết trước) và không nắm được ngữ cảnh xa trong câu
(phủ định, mệnh đề đối lập kiểu ``pin yếu nhưng màn hình đẹp''). Cài đặt:
\texttt{src/baseline/tfidf\_logreg.py}, chạy qua
\texttt{scripts/train\_baseline.py}.

\subsection{DistilBERT / BERT fine-tuned (domain laptop)}

Để khắc phục hạn chế về ngữ cảnh của baseline, nhóm fine-tune hai model
BERT-family theo kiểu \emph{sentence-pair classification} (BERT-SPC): input
là cặp \texttt{[CLS] sentence [SEP] aspect\_term [SEP]}, tokenize bằng
WordPiece, để model tự học mối liên hệ giữa câu và aspect term thông qua
cơ chế self-attention thay vì chỉ dựa vào tần suất từ như TF-IDF.

Nhóm fine-tune \texttt{distilbert-base-uncased} (lr=3e-5) và
\texttt{bert-base-uncased} (lr=2e-5), mỗi model chạy 3 seed (42/43/44)
$\times$ tối đa 20 epoch với early stopping (kiên nhẫn 3 epoch theo Macro-F1
trên tập valid) để tránh overfit khi budget epoch dư thừa. Hàm loss là
cross-entropy có trọng số lớp (tính theo tần suất nghịch đảo trên tập
train) để bù lệch giữa \texttt{positive}/\texttt{negative} và
\texttt{neutral}. Kết quả cuối được báo cáo dưới dạng trung bình $\pm$ độ
lệch chuẩn trên 3 seed, đo trên tập test giữ riêng.

Việc fine-tune chạy hoàn toàn trên Kaggle GPU miễn phí (không cần cài
transformers/torch ở máy local) -- notebook:
\texttt{notebooks/finetune\_\{distilbert,bert\}\_semeval\_laptop.ipynb}.

\subsection{Baseline tra cứu tần suất -- AsteLookupBaseline (domain nhà hàng)}

Baseline TF-IDF+LogReg ở trên không áp dụng được cho bài toán trích xuất bộ
ba vì nó cần biết trước aspect -- ASTE thì ngược lại, phải tự tìm aspect.
Nhóm xây một baseline phi-neural khác, phù hợp hơn: không cần train
gradient, chỉ ghi nhớ (a) tập các cụm aspect và (b) sentiment phổ biến nhất
của từng cụm opinion, cả hai đều thống kê trực tiếp từ tập train. Ở câu
test, thuật toán tra cứu nguyên văn các cụm đã thấy trong câu (khớp n-gram
dài nhất trước, tối đa 3 từ, quét trái sang phải không chồng lấn), rồi ghép
mỗi aspect tìm được với opinion \emph{gần nhất} trên câu (theo khoảng cách
token).

Baseline này có recall thấp theo thiết kế: bất kỳ cách diễn đạt aspect hay
opinion nào khác với những gì đã thấy trong tập train đều bị bỏ sót hoàn
toàn -- đây chính là ``sàn'' để đo mức độ cải thiện thực sự của một model
biết tổng quát hoá. Cài đặt: \texttt{src/baseline/aste\_lookup\_baseline.py},
chạy hoàn toàn local (không cần GPU) qua \texttt{scripts/eval\_aste\_baseline.py}.

\subsection{T5-small / T5-base / FLAN-T5-base fine-tuned cho ASTE (domain nhà hàng)}

Thay vì tagging từng token như repo nền, nhóm sinh trực tiếp bộ ba (aspect,
opinion, sentiment) bằng một mô hình text-to-text: câu thô được thêm một
prefix hướng dẫn cố định (\texttt{"extract aspect sentiment triplets: "})
làm input, target text có định dạng
\texttt{"aspect: X | opinion: Y | sentiment: Z ; ..."} (nối nhiều bộ ba
bằng dấu \texttt{;} nếu câu có nhiều aspect). Cách này vừa giải quyết được
bài toán trích xuất aspect (không cần gán nhãn trước), vừa cho thêm trường
\emph{opinion} làm ``lý do'' phục vụ bước sinh báo cáo -- điều mà pipeline
tagging kiểu Track~1 không có sẵn.

Ba checkpoint T5-family được fine-tune với \emph{cùng} hyperparameter (20
epoch, lr=3e-4, effective batch size 16 -- riêng T5-base dùng batch nhỏ hơn
kèm gradient accumulation lớn hơn để giữ effective batch không đổi do cần
nhiều bộ nhớ GPU hơn) để so sánh công bằng, chỉ đổi checkpoint:
\texttt{t5-small}, \texttt{t5-base}, \texttt{google/flan-t5-base}. Ở phần
dự đoán, chuỗi văn bản model sinh ra được parse bằng regex thành danh sách
bộ ba có cấu trúc.

Metric đánh giá là \emph{triplet-F1}: Precision/Recall/F1 theo phép giao
tập hợp (set overlap) giữa bộ ba dự đoán và bộ ba gốc (đã chuẩn hoá chữ
thường), micro-average trên toàn bộ tập test (không average theo từng câu
rồi lấy trung bình, để câu có nhiều bộ ba đóng góp đúng tỉ trọng). Hạ tầng:
ba notebook Kaggle độc lập, tự động tải dataset (clone GitHub nếu chưa có
sẵn trên Kaggle input) -- \texttt{notebooks/train-\{t5-small,t5-base,%
flan-t5-base\}-for-aste-*.ipynb}.

\subsection{Tổng hợp thống kê theo khía cạnh + top-10 lý do}

Ở cả hai track, bước tổng hợp gộp các cặp (aspect, sentiment) theo aspect,
đếm số lượt mỗi sentiment và lấy sentiment đa số. Với track nhà hàng (có
thêm trường opinion), bước này được mở rộng: trong mỗi nhóm
(aspect, sentiment), các cụm \emph{opinion} được xếp hạng theo tần suất, giữ
lại top-10 làm ví dụ ``lý do'' minh hoạ cho sentiment đó -- ví dụ aspect
\texttt{food} với sentiment \texttt{positive} có thể có lý do phổ biến nhất
là \texttt{great}, \texttt{good}, \texttt{delicious}.

Toàn bộ logic gộp số liệu này là thuần Python (không cần GPU), được viết và
test local (\texttt{src/report/aspect\_stats.py}, kiểm thử ở
\texttt{tests/test\_aspect\_stats.py}); chỉ riêng phần suy luận mô hình
(cần GPU) chạy trên Kaggle, sau đó logic tổng hợp đã test được inline
nguyên văn vào notebook để đảm bảo đúng kết quả như bản đã kiểm thử local.

\subsection{Sinh báo cáo bằng FLAN-T5}

\todo{Hoàng/Vinh/Hưng -- mô tả input representation (cách serialize
bảng thống kê thành prompt), cách dùng FLAN-T5 (few-shot/prompt-based
hay có fine-tune), và ví dụ input/output cụ thể. Tham khảo
\texttt{notebooks-output/report\_generation\_semeval\_laptop\_output.ipynb}
(branch \texttt{feat/report-generation-semeval-laptop-output}) và branch
\texttt{flan-t5-report-generation}.}

\subsection{Factual checker}

\todo{Hoàng/Vinh/Hưng -- mô tả cụ thể cách đối chiếu số liệu trong báo
cáo sinh ra với bảng thống kê gốc (rule-based, regex trích số liệu +
tên aspect, so khớp). Tham khảo \texttt{output/flan\_t5\_natural\_report.json}
làm ví dụ output thực tế (report + kết quả factual\_check).}


% =====================================================================
\section{Kết quả và đánh giá}
% =====================================================================

\subsection{Track 1 -- Phân loại sentiment theo khía cạnh (domain laptop)}

\begin{table}[h]
\centering
\begin{tabular}{lcc}
\toprule
Model & Accuracy & Macro-F1 \\
\midrule
Baseline TF-IDF+LogReg (3 lớp, cùng split BERT) & 0.6171 & 0.5581 \\
DistilBERT (3 seed) & $0.7447 \pm 0.0159$ & $0.6861 \pm 0.0235$ \\
\textbf{BERT-base (3 seed)} & $\mathbf{0.7627 \pm 0.0239}$ & $\mathbf{0.7123 \pm 0.0324}$ \\
\bottomrule
\end{tabular}
\caption{Track laptop: baseline vs.\ model cải tiến, đo trên cùng test
split (316 example, 3 lớp positive/negative/neutral).}
\label{tab:laptop-results}
\end{table}

\begin{figure}[h]
\centering
\begin{tikzpicture}
\begin{axis}[
    ybar, bar width=22pt,
    width=9.5cm, height=6.5cm,
    symbolic x coords={Baseline,DistilBERT,BERT-base},
    xtick=data,
    ymin=0, ymax=0.85,
    ylabel={Macro-F1},
    nodes near coords,
    nodes near coords align={vertical},
    every node near coord/.append style={font=\small},
    enlarge x limits=0.3,
]
\addplot[fill=blue!30] coordinates {(Baseline,0.5581) (DistilBERT,0.6861) (BERT-base,0.7123)};
\end{axis}
\end{tikzpicture}
\caption{Macro-F1 track laptop qua từng bước cải tiến.}
\label{fig:laptop-bar}
\end{figure}

BERT-base cải thiện +14.6 điểm Accuracy, +15.4 điểm Macro-F1 so với
baseline (đã chạy lại baseline trên đúng split 3 lớp để so sánh công
bằng -- xem thảo luận về nhãn \texttt{conflict} ở Mục~\ref{sec:findings}).
Lớp \texttt{neutral} là lớp yếu nhất ở cả hai mô hình.

\subsection{Track 2 -- Trích xuất bộ ba (domain nhà hàng)}

\begin{table}[h]
\centering
\begin{tabular}{lccc}
\toprule
Model & Triplet-F1 & Precision & Recall \\
\midrule
AsteLookupBaseline (phi-neural) & 0.3736 & 0.3222 & 0.4446 \\
T5-small & 0.7240 & 0.7238 & 0.7243 \\
\textbf{T5-base} & \textbf{0.7442} & 0.7609 & 0.7282 \\
FLAN-T5-base (lr=3e-4) & 0.5898 & 0.5910 & 0.5887 \\
FLAN-T5-base (lr=1e-4, thử lại) & 0.4159 & 0.4185 & 0.4133 \\
\bottomrule
\end{tabular}
\caption{Track nhà hàng: baseline vs.\ 3 model T5 so sánh, đo trên cùng
test split (1134 câu, 14res+15res+16res).}
\label{tab:aste-results}
\end{table}

\begin{figure}[h]
\centering
\begin{tikzpicture}
\begin{axis}[
    width=\textwidth, height=7.5cm,
    xlabel={Epoch}, ylabel={Dev triplet-F1},
    xmin=1, xmax=20, ymin=0.2, ymax=0.8,
    legend pos=south east,
    legend style={font=\footnotesize},
    grid=major,
]
\addplot[color=blue, mark=*, mark size=1pt] coordinates {
(1,0.5699) (2,0.6416) (3,0.6564) (4,0.6986) (5,0.6896) (6,0.7131)
(7,0.7201) (8,0.7188) (9,0.7345) (10,0.7293) (11,0.7379) (12,0.7409)
(13,0.7426) (14,0.7577) (15,0.7553) (16,0.7507) (17,0.7466) (18,0.7542)
(19,0.7489) (20,0.7525)
};
\addlegendentry{t5-small}
\addplot[color=red, mark=square*, mark size=1pt] coordinates {
(1,0.6773) (2,0.7207) (3,0.7404) (4,0.7433) (5,0.7570) (6,0.7736)
(7,0.7641) (8,0.7511) (9,0.7628) (10,0.7576) (11,0.7626) (12,0.7592)
(13,0.7666) (14,0.7692) (15,0.7607) (16,0.7612) (17,0.7644) (18,0.7683)
(19,0.7686) (20,0.7700)
};
\addlegendentry{t5-base}
\addplot[color=green!55!black, mark=triangle*, mark size=1.3pt] coordinates {
(1,0.2447) (2,0.2526) (3,0.3114) (4,0.3944) (5,0.4375) (6,0.4907)
(7,0.5022) (8,0.5397) (9,0.5722) (10,0.5858) (11,0.6177) (12,0.6320)
(13,0.6499) (14,0.6628) (15,0.6660) (16,0.6771) (17,0.6881) (18,0.6820)
(19,0.6795) (20,0.6801)
};
\addlegendentry{flan-t5-base (lr=3e-4)}
\addplot[color=orange!90!black, mark=diamond*, mark size=1.3pt] coordinates {
(1,0.2436) (2,0.2977) (3,0.3027) (4,0.2980) (5,0.2915) (6,0.3322)
(7,0.3359) (8,0.3542) (9,0.4046) (10,0.4093) (11,0.4168) (12,0.4228)
(13,0.4561) (14,0.4505) (15,0.4594) (16,0.4701) (17,0.4683) (18,0.4742)
(19,0.4734) (20,0.4703)
};
\addlegendentry{flan-t5-base (lr=1e-4)}
\end{axis}
\end{tikzpicture}
\caption{Dev triplet-F1 theo epoch, cả 4 lần chạy (số liệu lấy trực tiếp từ
log training trong \texttt{notebooks-output/}). t5-small và t5-base hội tụ
và bão hoà rõ trong khoảng epoch 10--15; cả hai đường flan-t5-base vẫn đang
tăng dần đều tới hết epoch 20, chưa hề bão hoà -- minh chứng trực quan cho
phát hiện ở Mục~\ref{sec:findings}: vấn đề là cần nhiều epoch hơn, không
phải learning rate.}
\label{fig:triplet-f1-curves}
\end{figure}

T5-base cải thiện gần gấp đôi triplet-F1 so với baseline tra cứu
(0.3736 $\to$ 0.7442) và được chọn làm model chính cho track này.
Hình~\ref{fig:triplet-f1-curves} cho thấy rõ lý do FLAN-T5-base thua kém:
không phải nó ``kém hơn'' về bản chất mà đơn giản là chưa hội tụ trong ngân
sách 20 epoch đã cấp cho cả 4 lần chạy. Xem Mục~\ref{sec:findings} để phân
tích sâu hơn.

\subsection{Chất lượng bảng tổng hợp aspect + lý do}

\begin{table}[h]
\centering
\begin{tabular}{lcc}
\toprule
& Track laptop (BERT) & Track nhà hàng (T5-base) \\
\midrule
Majority-sentiment agreement & 93.28\% & 97.25\% \\
(gold vs.\ predicted, cùng aspect) & (236/253) & (531/546) \\
\bottomrule
\end{tabular}
\caption{Độ khớp giữa kết luận sentiment đa số (gold) và kết luận từ
model (predicted) trên cùng một aspect.}
\label{tab:agreement}
\end{table}

Track nhà hàng có agreement cao hơn hẳn -- model dự đoán đồng thời cả
aspect/opinion/sentiment thay vì chỉ suy luận sentiment cho aspect đã
biết trước.

\subsection{Demo quy mô lớn -- Yelp Restaurant Reviews}

Thay vì Amazon Reviews (dự kiến ban đầu, category Electronics -- lệch
domain hoàn toàn với model T5 train trên domain nhà hàng), nhóm đổi
sang \href{https://www.kaggle.com/datasets/farukalam/yelp-restaurant-reviews}{Yelp
Restaurant Reviews} để domain khớp với model, không có gold label (chỉ
demo quy mô lớn, không benchmark).

Chạy trên 2000 review mẫu (seed=42) $\to$ 14285 câu (sau khi tách câu
từ review nhiều câu) $\to$ 14700 bộ ba $\to$ 1224 aspect ($\geq 2$ lượt
nhắc). Mẫu lấy được nghiêng nhiều về nhóm quán tráng miệng/bakery: top
aspect \texttt{ice cream} (765 lượt, positive), \texttt{place} (633),
\texttt{flavors} (237), cùng \texttt{staff}/\texttt{service}/
\texttt{donuts}/\texttt{pastries}/\texttt{bakery} -- kết quả hợp lý,
đúng domain.

\subsection{Đánh giá factual checker}

\todo{Hoàng/Vinh/Hưng -- báo cáo \% số liệu/aspect trong report khớp
với bảng thống kê gốc, trên bao nhiêu report được sinh ra, có ví dụ cụ
thể (report bị factual checker gắn cờ + lý do).}

\subsection{Error analysis}

\todo{Chưa làm (Tuần 5). Gợi ý phân tích: (1) so sánh lỗi phổ biến của
BERT trên lớp \texttt{neutral} (F1 thấp nhất) -- loại câu nào hay bị
nhầm; (2) các bộ ba T5-base dự đoán sai trên tập test (dùng
\texttt{test\_predictions.csv} lưu trong mỗi model checkpoint); (3) các
aspect noise trong demo Yelp (vd.\ \texttt{extract} -- 196 lượt, nhiều
khả năng model trích cụt từ ``vanilla extract'').}


% =====================================================================
\section{Khám phá mới / Thảo luận}
\label{sec:findings}
% =====================================================================

\begin{enumerate}
    \item \textbf{Macro-F1 rất nhạy với một lớp thiểu số học kém.} Baseline
        đo trên bản dataset 4 lớp cũ (có thêm \texttt{conflict}, chỉ
        chiếm $<0.3\%$ dữ liệu) cho Macro-F1 chỉ 0.4266 -- thấp hơn hẳn
        so với 0.5581 khi đo trên đúng 3 lớp hiện tại, dù Accuracy gần
        như không đổi (0.6211 vs.\ 0.6171). Nguyên nhân: baseline không
        bao giờ đoán đúng lớp \texttt{conflict}, kéo trung bình không
        trọng số (macro) xuống rất nặng dù lớp đó gần như không ảnh
        hưởng accuracy. Bài học: khi so sánh hai mô hình đo trên hai
        schema nhãn khác nhau, Macro-F1 cần được đối chiếu lại trên
        cùng schema mới kết luận công bằng được.

    \item \textbf{Đừng vội kết luận nguyên nhân từ một hyperparameter.}
        FLAN-T5-base thua kém rõ rệt so với T5-base cùng kích thước khi
        train với cùng lr=3e-4 (0.5898 vs.\ 0.7442 triplet-F1). Giả
        thuyết ban đầu: lr quá cao khiến loss ban đầu bất ổn (xấp xỉ
        22, so với hai model kia ổn định ngay từ đầu). Thử lại với
        lr=1e-4 (giảm 10 lần) để kiểm chứng -- nhưng kết quả \emph{tệ
        hơn} (0.4159), và Hình~\ref{fig:triplet-f1-curves} cho thấy
        \texttt{eval\_triplet\_f1} của cả hai lần chạy flan-t5-base vẫn
        tăng dần đều tới hết epoch 20, chưa hề bão hoà -- trong khi
        t5-small/t5-base đã bão hoà rõ từ epoch 10--15. Kết luận đúng:
        lr thấp hơn chỉ làm hội tụ chậm hơn, không giải quyết được gì --
        nguyên nhân thực sự nhiều khả năng là FLAN-T5-base (đã
        instruction-tuned trên nhiều task đa dạng) cần nhiều epoch hơn
        hẳn để thích nghi lại với định dạng trích xuất terse này, không
        phải vấn đề learning rate.

    \item \textbf{Sinh trực tiếp bộ ba giải quyết luôn bài toán trích
        xuất aspect, và cho kết quả tổng hợp đáng tin hơn.} So với
        pipeline cần biết trước aspect (track laptop), pipeline T5-ASTE
        (track nhà hàng) không chỉ tự tìm được aspect mà còn cho
        majority-sentiment agreement cao hơn hẳn (97.25\% vs.\ 93.28\%)
        khi tổng hợp thống kê -- gợi ý rằng học đồng thời cả ba thông
        tin (aspect, opinion, sentiment) trong một lần sinh giúp mô
        hình nhất quán hơn so với suy luận sentiment riêng lẻ cho từng
        aspect.
\end{enumerate}


% =====================================================================
\section{Hạn chế}
% =====================================================================

\begin{itemize}
    \item \textbf{AsteLookupBaseline giới hạn từ vựng nghiêm ngặt}: chỉ
        khớp nguyên văn cụm từ đã thấy trong train, không có bất kỳ hình
        thức chuẩn hoá ngữ nghĩa nào (đồng nghĩa, biến thể hình thái).
        Recall 0.4446 phần lớn phản ánh giới hạn này hơn là độ khó thực sự
        của bài toán.
    \item \textbf{Không dùng POS tagging thật} cho baseline ASTE (chỉ tra
        cứu + ghép theo khoảng cách token) -- một baseline rule-based
        ``chuẩn'' hơn trong literature thường dùng POS tagger để tìm ứng
        viên danh từ/tính từ trước khi ghép cặp; nhóm bỏ qua để tránh thêm
        dependency (NLTK/spaCy) không cần thiết cho một baseline tham
        chiếu.
    \item \textbf{Track nhà hàng và track laptop không dùng chung domain}
        -- kết quả hai track không so sánh chéo được (khác dataset, khác
        task). Việc BERT (track laptop) và T5-base (track nhà hàng) đều
        ``tốt hơn baseline'' không có nghĩa hai track tương đương nhau về
        độ khó.
    \item \textbf{Demo Yelp không có gold label} -- độ tin cậy của bảng
        thống kê trên dữ liệu này chỉ được suy luận gián tiếp từ agreement
        97.25\% đo trên dữ liệu có nhãn cùng domain (SemEval Restaurant),
        không phải đo trực tiếp trên chính dữ liệu Yelp.
    \item \textbf{Factual checker rule-based, không xử lý được diễn đạt
        tự do}: cách đối chiếu bằng regex (theo kế hoạch ở
        Mục~\ref{sec:methods}) sẽ bỏ sót các trường hợp báo cáo diễn đạt
        đúng ý nhưng không dùng đúng định dạng số liệu kỳ vọng.
        \todo{Hoàng/Vinh/Hưng -- xác nhận/bổ sung hạn chế thực tế sau khi
        có kết quả đánh giá.}
\end{itemize}


% =====================================================================
\section{Kết luận}
% =====================================================================

Nhóm đã xây dựng và so sánh đầy đủ pipeline baseline--đến--model cải tiến
trên hai track: (1) phân loại sentiment theo khía cạnh đã biết trước, domain
laptop (TF-IDF+LogReg $\to$ DistilBERT $\to$ BERT-base, cải thiện Macro-F1
+15.4 điểm), và (2) trích xuất trực tiếp bộ ba khía cạnh--ý kiến--sentiment
từ câu thô, domain nhà hàng (AsteLookupBaseline $\to$ T5-small $\to$
T5-base, cải thiện triplet-F1 gần gấp đôi). Hướng thứ hai là một bổ sung
quan trọng: vừa giải quyết được bài toán trích xuất aspect mà track đầu
không làm được, vừa cho thêm trường \emph{opinion} làm ``lý do'' phục vụ
bước sinh báo cáo, vừa cho độ nhất quán tổng hợp cao hơn hẳn (97.25\% so
với 93.28\%). Nhóm khuyến nghị dùng T5-base (track nhà hàng) làm nguồn dữ
liệu chính cho bước sinh báo cáo cuối cùng.

\todo{Bổ sung sau khi có đủ kết quả factual checker + error analysis: (1)
tỉ lệ báo cáo bị factual checker gắn cờ và loại lỗi phổ biến nhất; (2)
tổng kết ngắn gọn insight từ error analysis; (3) hướng phát triển tiếp
nếu có thêm thời gian (vd.\ mở rộng track nhà hàng sang aspect-category
cố định, hoặc thử fine-tune FLAN-T5-base với epoch budget lớn hơn để
kiểm chứng giả thuyết ở Mục~\ref{sec:findings}).}


% =====================================================================
\section*{Phân công nhóm}
% =====================================================================

\begin{table}[h]
\centering
\begin{tabular}{ll}
\toprule
Việc & Thành viên \\
\midrule
Data pipeline + baseline (Tuần 3) & Đạt \\
Fine-tune BERT/T5 + tổng hợp thống kê (Tuần 4) & Đạt, Sơn \\
FLAN-T5 sinh report + factual checker (Tuần 4) & Hoàng, Vinh, Hưng \\
Đánh giá \& phân tích (Tuần 5) & \todo{phân công} \\
Viết báo cáo + slides (Tuần 6) & \todo{phân công} \\
\bottomrule
\end{tabular}
\end{table}


% =====================================================================
\appendix
\section{Phụ lục}
% =====================================================================

\subsection{Bảng hyperparameter đầy đủ}

\begin{table}[h]
\centering
\small
\begin{tabular}{lccccc}
\toprule
Model & Epoch & LR & Batch (effective) & Seed & Early stopping \\
\midrule
DistilBERT & $\leq$20 & 3e-5 & 16 & 42/43/44 & patience=3 (Macro-F1) \\
BERT-base & $\leq$20 & 2e-5 & 16 & 42/43/44 & patience=3 (Macro-F1) \\
T5-small (ASTE) & 20 & 3e-4 & 16 & 42 & không \\
T5-base (ASTE) & 20 & 3e-4 & 16 & 42 & không \\
FLAN-T5-base (ASTE) & 20 & 3e-4 rồi 1e-4 & 16 & 42 & không \\
\bottomrule
\end{tabular}
\caption{Hyperparameter đầy đủ cho các model fine-tune. Chi tiết từng
notebook nằm ở \texttt{notebooks/}.}
\end{table}

\subsection{Ví dụ output thực tế -- Track nhà hàng}

Ví dụ một hàng trong bảng tổng hợp (\texttt{output/aspect\_reasons\_%
restaurant.json}), aspect \texttt{food} (827 lượt nhắc, T5-base, domain
nhà hàng):

\begin{quote}
\texttt{positive} (621 lượt): lý do phổ biến nhất --
\texttt{great}(109), \texttt{good}(100), \texttt{delicious}(39),
\texttt{excellent}(37). \\
\texttt{negative} (164 lượt): lý do phổ biến nhất --
\texttt{mediocre}(10), \texttt{bad}(7), \texttt{overpriced}(5).
\end{quote}

\subsection{Ví dụ output thực tế -- Track laptop / báo cáo sinh tự động}

\todo{Hoàng/Vinh/Hưng -- dán 1--2 ví dụ report thật (input bảng thống kê
+ report sinh ra + kết quả factual\_check) từ
\texttt{output/flan\_t5\_natural\_report.json} vào đây.}


% =====================================================================
\begin{thebibliography}{99}

\bibitem{semeval2014} Pontiki, M., Galanis, D., Pavlopoulos, J., Papageorgiou, H., Androutsopoulos, I., \& Manandhar, S. (2014). \emph{SemEval-2014 Task 4: Aspect Based Sentiment Analysis}. \url{https://alt.qcri.org/semeval2014/task4/}

\bibitem{berte2eabsa} Li, X., Bing, L., Zhang, W., \& Lam, W. (2019). \emph{Exploiting BERT for End-to-End Aspect-Based Sentiment Analysis}. \url{https://github.com/lixin4ever/BERT-E2E-ABSA}

\bibitem{peng2020aste} Peng, H., Xu, L., Bing, L., Huang, F., Lu, W., \& Si, L. (2020). \emph{Knowing What, How and Why: A Near Complete Solution for Aspect-Based Sentiment Analysis}. AAAI 2020.

\bibitem{xu2020aste} Xu, L., Chia, Y. K., \& Bing, L. (2020). \emph{Position-Aware Tagging for Aspect Sentiment Triplet Extraction}. EMNLP 2020.

\bibitem{t5} Raffel, C., et al. (2020). \emph{Exploring the Limits of Transfer Learning with a Unified Text-to-Text Transformer}. Journal of Machine Learning Research.

\bibitem{flant5} Chung, H. W., et al. (2022). \emph{Scaling Instruction-Finetuned Language Models}. arXiv:2210.11416.

\bibitem{t5totto} Kale, M., \& Rastogi, A. (2020). \emph{Text-to-Text Pre-Training for Data-to-Text Tasks}. INLG 2020. \url{https://aclanthology.org/2020.inlg-1.14.pdf}

\bibitem{factcc} Kryscinski, W., McCann, B., Xiong, C., \& Socher, R. (2020). \emph{Evaluating the Factual Consistency of Abstractive Text Summarization} (FactCC).

\bibitem{consistency2024} \emph{Consistency Is the Key: Improving Factuality via Consistency-based Checking}. arXiv:2511.12236.

\bibitem{devlin2019bert} Devlin, J., Chang, M.-W., Lee, K., \& Toutanova, K. (2019). \emph{BERT: Pre-training of Deep Bidirectional Transformers for Language Understanding}. NAACL 2019.

\bibitem{sanh2019distilbert} Sanh, V., Debut, L., Chaumond, J., \& Wolf, T. (2019). \emph{DistilBERT, a distilled version of BERT: smaller, faster, cheaper and lighter}. arXiv:1910.01108.

\end{thebibliography}

\end{document}
